\documentclass[12pt]{book}
\usepackage{fontspec}
\setmainfont{Noto Serif CJK JP}
\usepackage[a4paper,textwidth=13cm,top=2.5cm,bottom=2.5cm]{geometry}
\usepackage{titlesec}
\renewcommand{\contentsname}{Inhaltsverzeichnis}
\renewcommand{\chaptername}{Kapitel}
\emergencystretch=3em
\tolerance=2000

\title{Warum に kein Dativ ist\\\large Japanische Grammatik in Bildern}
\author{Markus Bauernfeind}
\date{2026}

\begin{document}
\pagestyle{plain}
\maketitle
\tableofcontents

\chapter{Vorwort}

Die meisten Lehrbücher erklären japanische Grammatik, indem sie diese in deutsche Kategorien pressen: Partikeln werden zu Fällen erklärt, Verbformen zu Zeiten, Konstruktionen zu Übersetzungsmustern. Das funktioniert eine Weile — bis es nicht mehr funktioniert, weil Japanisch nach einer ganz anderen inneren Logik gebaut ist als Deutsch.

Dieses Buch versucht etwas anderes: zu zeigen, welches Bild ein japanischer Muttersprachler im Kopf hat, wenn er sich für eine Partikel oder eine Form entscheidet. に ist kein "Dativ" — es ist ein Bezugspunkt, an den sich etwas andockt. を ist kein "Akkusativ" — es ist eine Ausdehnung, die vollständig durchmessen wird. Wer dieses innere Bild versteht, muss sich nicht mehr merken, welche Partikel "man in diesem Fall benutzt" — er erkennt, warum sie die naheliegende Wahl ist.

\section*{Was dieses Buch nicht ist}

Es ist keine linguistische Abhandlung mit Vollständigkeitsanspruch. Sprachliche Phänomene sind komplex, und ein Lernbild muss nicht jede akademische Randdebatte abbilden. Die hier vorgestellten Bilder erklären die häufigsten und lehrreichsten Fälle einer Partikel oder Form — nicht notwendigerweise jede einzelne Randverwendung.

Es ist auch keine Sammlung endgültiger, "bewiesener" Wahrheiten. Selbst Fachlinguisten sind sich bei manchen Partikeln — etwa に oder は — nicht vollständig einig, wie ihre Bedeutung im Kern zu fassen ist. Die Bilder in diesem Buch sind funktionierende, sorgfältig geprüfte Lernvereinfachungen, keine behaupteten letzten Wahrheiten. Wo ein Kapitel eine Grenze oder einen offenen Punkt kennt, wird das offen benannt statt verschwiegen.

\section*{Wie die Kapitel aufgebaut sind}

Jedes Kapitel widmet sich einer Partikel oder Form und stellt das Bild in den Mittelpunkt: die Vorstellung, die einem Sprecher naheliegt, wenn er sich für diese Form entscheidet. Beispiele veranschaulichen das Bild an konkreten Sätzen. Wo es einen echten, lehrreichen Grenzfall gibt, wird er gezeigt — nicht um zu verunsichern, sondern weil solche Fälle oft am deutlichsten zeigen, wie das Bild funktioniert. Wo die Herkunft eines Wortes zum Verständnis beiträgt, steht sie als kurze Randnotiz am Ende des Kapitels.

Nicht jedes Kapitel braucht jeden dieser Bausteine. Manche Partikeln haben einen interessanten Grenzfall, andere nicht; manche eine aufschlussreiche Etymologie, andere nicht. Ein Kapitel wird nicht künstlich aufgefüllt, nur um einem Schema zu folgen.

\section*{Für wen dieses Buch gedacht ist}

Für Menschen, die Japanisch lernen — nicht für ein linguistisches Fachpublikum. Ob als Ergänzung zu einem Kurs, als Nachschlagewerk oder zum Selbststudium: das Ziel ist immer dasselbe. Ein Gefühl dafür zu entwickeln, warum die japanische Sprache so funktioniert, wie sie funktioniert — nicht auswendig zu lernen, sondern zu verstehen.

\section*{Hinweis zur Entstehung}

Bei der Erarbeitung dieses Buches wurde Claude (Anthropic) als KI-Assistent eingesetzt — für Recherche, sprachliche Ausarbeitung und Gegenbeispielprüfung der vorgestellten Bilder. Alle Konzepte, Entscheidungen und die inhaltliche Beurteilung stammen vom Autor.


\chapter{に — Die Umfeld-Partikel}

\section*{Das Bild}

Stell dir に als einen Finger vor, der auf einen \textbf{festen Punkt} zeigt — einen Ort, einen Zeitpunkt, ein Ziel, einen Zustand. に markiert nicht \emph{Bewegung durch etwas}, sondern \emph{das, worauf sich etwas ausrichtet oder worin etwas verankert ist}.

Wenn ein Japaner に sagt, denkt er nicht in der Kategorie "Dativ" oder "Ort" oder "Zeit" — das sind deutsche Schubladen, die nachträglich über eine einzige, einheitliche Vorstellung gelegt werden. Im Kopf des Sprechers gibt es nur \emph{ein} Bild: \textbf{etwas wird an einem bestimmten Punkt verankert.}

\section*{Die wichtigsten Fälle}

\textbf{1. Ort der Existenz (nicht der Handlung)}
\par\quad 東京に住んでいます。\\
\par\quad Ich wohne in Tokyo.\\

Der Punkt, an dem etwas \emph{ist}, nicht durch den es sich bewegt. Vergleiche mit で (Ort der Handlung) — dort geht es um die Bühne, hier um die Verankerung.

\textbf{2. Zeitpunkt}
\par\quad 7時に起きます。\\
\par\quad Ich stehe um 7 Uhr auf.\\

Zeit wird genauso behandelt wie Ort: ein fester Punkt, an dem etwas verankert ist.

\textbf{3. Ziel/Richtung einer Bewegung}
\par\quad 学校に行きます。\\
\par\quad Ich gehe zur Schule.\\

Der Zielpunkt der Bewegung — nicht der Weg dorthin (das wäre eher を bei bestimmten Verben), sondern der Punkt, auf den die Bewegung ausgerichtet ist.

\textbf{4. Empfänger/Richtung einer Handlung}
\par\quad 友達に手紙を書きます。\\
\par\quad Ich schreibe meinem Freund einen Brief.\\

Auch die Handlung ist hier auf einen Punkt (die Person) ausgerichtet — derselbe Mechanismus wie bei Ort und Zeit, nur mit einer Person als Anker.

\textbf{5. Zustandsänderung (Ziel eines Wandels)}
\par\quad 先生になりたいです。\\
\par\quad Ich möchte Lehrer werden.\\

Der Zustand, in dem man "ankommt" — die Zustandsänderung hat einen Zielpunkt, genau wie eine Bewegung.

\section*{Ein interessanter Grenzfall}

Weil に selbst \textbf{richtungsneutral} ist (die genaue Richtung kommt vom Verb, nicht von に selbst), kann derselbe Satz manchmal zwei Lesarten haben.

\par\quad 議長に推薦された。\\

Das kann heißen "ich wurde \emph{zum} Vorsitzenden vorgeschlagen" (に als Ziel) — oder "ich wurde \emph{vom} Vorsitzenden vorgeschlagen" (に als Ursprung/Urheber der Handlung). Beides ist grammatisch korrekt, der Kontext muss entscheiden. Das ist kein Fehler im Bild, sondern eine direkte Folge davon: wenn に nur den Bezugspunkt markiert und nicht die Richtung selbst, dann \emph{muss} echte Mehrdeutigkeit entstehen, sobald der Kontext keine eindeutige Richtung vorgibt.

\bigskip\hrule\bigskip

\emph{Für Interessierte:} に geht auf ein altjapanisches Element zurück, das ursprünglich einen statischen Ort bezeichnete — im Gegensatz zu へ, das historisch eine Richtung/Bewegung hin zu ausdrückte. Diese Wurzel im "Verankerungspunkt" erklärt, warum sich aus に so unterschiedlich wirkende Verwendungen entwickeln konnten.



\chapter{で — Die Grenz-Partikel}

\section*{Das Bild}

Stell dir で als eine gezogene Grenze vor — ein Kreis um einen Bereich, ein Strich als Limit, ein Punkt, an dem ein Zustand in einen anderen kippt. で sagt: \emph{hier ist die Grenze, innerhalb derer, bis zu der, oder an der etwas geschieht.}

Anders als に (ein Punkt, an dem etwas verankert ist) markiert で immer eine \textbf{Abgrenzung} — mal als Fläche (ein Bereich, in dem etwas stattfindet), mal als Limit (eine Obergrenze), mal als Scheidepunkt (der Übergang von einem Zustand zum nächsten).

\section*{Die wichtigsten Fälle}

\textbf{1. Ort der Handlung}
\par\quad 公園で遊びます。\\
\par\quad Ich spiele im Park.\\

Der Park ist der abgegrenzte Bereich, innerhalb dessen die Handlung stattfindet — nicht der Ort, an dem etwas einfach \emph{existiert} (das wäre に), sondern die Bühne, auf der etwas \emph{passiert}.

\textbf{2. Mittel/Werkzeug}
\par\quad ペンで書きます。\\
\par\quad Ich schreibe mit einem Stift.\\

Der Stift steckt die Grenze ab, innerhalb derer das Schreiben möglich wird — das Mittel als abgegrenzter Ermöglichungsraum der Handlung.

\textbf{3. Material}
\par\quad 紙で作ります。\\
\par\quad Ich mache es aus Papier.\\

Die Ausgangsgrenze, aus der heraus etwas entsteht.

\textbf{4. Summe oder Frist}
\par\quad 1000円で買いました。\\
\par\quad Ich habe es für 1000 Yen gekauft.\\

\par\quad 一週間でできます。\\
\par\quad Es ist in einer Woche fertig.\\

Eine Obergrenze — beim Geld ein Limit, bei der Zeit eine Frist.

\textbf{5. Ursache}
\par\quad 地震で電車が止まりました。\\
\par\quad Wegen des Erdbebens hielten die Züge an.\\

Hier ist die Grenze ein \textbf{Scheidepunkt}: der Moment, an dem ein Zustand (Züge fahren) in einen anderen (Züge stehen) kippt. Dasselbe Prinzip trägt auch あとで ("danach", Scheidepunkt zwischen Abgeschlossenem und Beginnendem) und begründendes ので.

\textbf{6. Vergleichsbereich}
\par\quad クラスの中で一番背が高いです。\\
\par\quad In der Klasse ist er der Größte.\\

中で ("innerhalb der Mitte") grenzt explizit die Gruppe ab, innerhalb derer der Vergleich gilt — dieselbe Grenzsetzungs-Funktion wie bei den anderen Fällen, nur besonders sichtbar.

\bigskip\hrule\bigskip

\emph{Für Interessierte:} で geht auf にて zurück, eine ältere Verbindung aus に und dem Partikel て. Die heutige Bandbreite der Verwendungen — Ort, Mittel, Material, Limit, Ursache — lässt sich einheitlich auf den Gedanken der gesetzten Grenze zurückführen, die den Rahmen einer Handlung absteckt.


\chapter{へ — Die Richtungs-Partikel}

\section*{Das Bild}

Stell dir へ als einen Pfeil vor, der sich noch \textbf{auf dem Weg} befindet — nicht der Punkt, an dem man ankommt, sondern die Bewegung selbst, die dorthin führt. Während に den Ankunftspunkt markiert (das Ergebnis), markiert へ den Prozess, den Verlauf, das Sich-dorthin-Bewegen.

Diese beiden Partikel überschneiden sich oft im Alltag — 東京に行く und 東京へ行く sagen beide "ich gehe nach Tokyo". Der feine Unterschied bleibt aber bestehen: に betont, \emph{wo} man ankommt, へ betont, \emph{dass man sich dorthin bewegt}.

\section*{Die wichtigsten Fälle}

\textbf{1. Bewegungsrichtung}
\par\quad 東京へ行きます。\\
\par\quad Ich fahre nach Tokyo (in Richtung Tokyo).\\

Der Fokus liegt auf der Bewegung selbst, weniger auf dem exakten Zielpunkt.

\textbf{2. Abstrakte Entwicklung}
\par\quad 夏へと向かっています。\\
\par\quad Es geht auf den Sommer zu.\\

Auch ohne physische Bewegung trägt das Bild: ein Prozess, eine Entwicklung, die sich auf etwas hin bewegt — "ein Schritt weiter" (一歩先へ) funktioniert genauso.

\textbf{3. Adressat einer Handlung}
\par\quad 恋人へ手紙を書きます。\\
\par\quad Ich schreibe der Geliebten einen Brief.\\

Hier markiert へ nicht nur den Empfänger, sondern betont den \emph{Verlauf} der Handlung — der Brief ist unterwegs, der Akt des Schreibens ist auf die Person hin gerichtet, nicht nur ein Endpunkt-Kontakt.

\section*{Ein interessanter Grenzfall}

Bei direktem, unmittelbarem Körperkontakt wirkt へ unnatürlich:

\par\quad 猫が私に噛みつきました。\\
\par\quad Die Katze hat mich gebissen.\\

Hier nur に, nicht へ. Das passt genau ins Bild: ein Biss ist kein Prozess, der sich über eine Strecke entwickelt — er ist ein diskreter Endpunkt-Kontakt, sofort da. Wo kein Verlauf, keine Bewegung-hin-zu stattfindet, hat へ nichts zu markieren.

\bigskip\hrule\bigskip

\emph{Für Interessierte:} へ geht auf das Nomen 辺 zurück ("Umgebung, Gegend" — wie in 道の辺, "Wegesrand"). Aus diesem räumlichen Flächenbegriff entwickelte sich im Lauf der Sprachgeschichte die Richtungspartikel — und drang später zunehmend auch in die Domäne von に (den Ankunftspunkt) ein, was die heutige Überlappung beider Partikel erklärt.


\chapter{を — Die Durchgangs-Partikel}

\section*{Das Bild}

Stell dir を als eine Strecke vor, die vollständig durchmessen, bearbeitet oder verlassen wird. を markiert nicht einfach "das Objekt" im Sinne einer grammatischen Kategorie — es markiert das, was die Handlung \textbf{ganz durchläuft, ganz erfasst oder ganz hinter sich lässt.}

Ein Objekt, das mit を markiert wird, wird von der Handlung vollständig durchdrungen oder verändert. Ein Ort, der mit を markiert wird, wird von der Bewegung vollständig durchquert. Ein Ausgangspunkt, der mit を markiert wird, wird von der Handlung vollständig verlassen. Immer derselbe Gedanke: die volle Ausdehnung von etwas wird durchmessen.

\section*{Die wichtigsten Fälle}

\textbf{1. Objekt einer Handlung}
\par\quad 小説を書きます。\\
\par\quad Ich schreibe einen Roman.\\

Der Roman entsteht vollständig durch die Handlung — von を erfasst als vollständig durchdrungener Wirkungsbereich.

\textbf{2. Durchquerter Ort}
\par\quad 山道を行きます。\\
\par\quad Ich gehe den Bergpfad entlang.\\

Hier ist を der Weg selbst, den die Bewegung vollständig durchmisst — nicht das Ziel (das wäre に oder へ), sondern die Strecke.

\textbf{3. Trennungs-/Ausgangspunkt}
\par\quad 大学を卒業します。\\
\par\quad Ich schließe die Universität ab (verlasse sie).\\

Die Universität wird als vollständig durchlaufene Phase hinter sich gelassen — derselbe Durchgangs-Gedanke, nur am Ende statt in der Mitte der Bewegung.

\textbf{4. Zeitliche Ausdehnung}
\par\quad 長い年月を過ごします。\\
\par\quad Ich verbringe lange Jahre.\\

Auch Zeit lässt sich als Strecke denken, die durchmessen wird.

\section*{Ein scharfer Kontrast zu で}

海を泳ぐ und 海で泳ぐ sehen fast gleich aus, bedeuten aber Verschiedenes:

\par\quad 海を泳ぎます。\\
\par\quad Ich durchschwimme das Meer.\\

\par\quad 海で泳ぎます。\\
\par\quad Ich schwimme im Meer.\\

Mit を wird das gesamte Meer als zu bewältigende Strecke gedacht — man durchmisst es vollständig. Mit で (der Grenz-Partikel) ist das Meer nur der Bereich, \emph{innerhalb} dessen lokal geschwommen wird, ohne Anspruch auf vollständige Durchquerung. Zwei Partikel, zwei völlig unterschiedliche Bilder — beide bleiben ihrem eigenen Bild treu.

\section*{Ein zweiter Kontrast: を und に bei Zwang und Erlaubnis}

Bei manchen kausativen Sätzen (jemanden zu etwas veranlassen) gibt es eine Wahl zwischen を und に — und der Unterschied ist genau der, den man erwarten würde:

\par\quad 子供を勉強させます。 → Zwang: das Kind wird vollständig durch das Lernen "hindurchgezwungen".\\
\par\quad 子供に勉強させます。 → Erlaubnis/Duldung: das Kind darf lernen, に markiert nur den lockeren Bezugspunkt.\\

を zwingt vollständig hindurch, に lässt nur zu. Zwei unabhängige Partikelbilder, die hier zusammen einen echten Bedeutungsunterschied erklären.


\chapter{が — Die Identifikations-Partikel}

\section*{Das Bild}

Stell dir が als einen Finger vor, der zum ersten Mal auf etwas zeigt: "Das hier!" が identifiziert eine Entität, die dem Zuhörer noch nicht als Thema bekannt ist, und macht sie zum direkten Träger dessen, was im Satz ausgesagt wird.

Während は (die Themapartikel) etwas Bekanntes als Rahmen für die Aussage setzt, bringt が etwas Neues, noch nicht Etabliertes ins Spiel und sagt: \emph{das} ist es, worüber jetzt gesprochen wird.

\section*{Die wichtigsten Fälle}

\textbf{1. Subjekt einer neu wahrgenommenen Tatsache}
\par\quad 雨が降っています。\\
\par\quad Es regnet.\\

Der Regen wird gerade als Tatsache erfasst und benannt — keine Vorerwähnung nötig.

\textbf{2. Objekt bei Zustands-, Gefühls- und Wahrnehmungsverben}
\par\quad サッカーが好きです。\\
\par\quad Ich mag Fußball.\\

\par\quad 音楽が聞こえます。\\
\par\quad Ich höre Musik.\\

Auch hier identifiziert が den Träger des Zustands oder der Wahrnehmung — grammatisch "Objekt", aber demselben Identifikations-Gedanken folgend.

\textbf{3. Existenz}
\par\quad 猫がいます。\\
\par\quad Da ist eine Katze.\\

が identifiziert, was existiert.

\textbf{4. Exklusive Identifikation}
\par\quad あいつが犯人です。\\
\par\quad \emph{Er} ist der Täter (und kein anderer).\\

Hier wird が besonders scharf: nicht nur "eine neue Tatsache", sondern "genau dieser und kein anderer" — eine stärker betonte Spielart derselben Grundfunktion.

\textbf{5. Fragewörter}
\par\quad 誰が来ますか。\\
\par\quad Wer kommt?\\

Fragewörter stehen praktisch immer mit が — folgerichtig, denn ein Fragewort ist per Definition unbekannt und kann daher nicht als bereits etablierter Ausgangspunkt einer Aussage dienen.


\chapter{は — Die Herausgreifungs-Partikel}

\section*{Das Bild}

Stell dir は als eine Hand vor, die ein Element aus dem Satz herausnimmt und in die Mitte stellt: "Darüber reden wir jetzt." は ist keine gewöhnliche Kasuspartikel wie が oder を — sie ersetzt diese sogar oft. Stattdessen \textbf{hebt} は ein Element heraus und macht es zum Ausgangspunkt für alles Folgende.

Diese Herausgreifung hat zwei Gesichter, die auf denselben Mechanismus zurückgehen:

\section*{Die wichtigsten Fälle}

\textbf{1. Thema}
\par\quad 象は鼻が長いです。\\
\par\quad Was Elefanten betrifft — der Rüssel ist lang.\\

は setzt den Rahmen: alles Folgende wird im Bezug auf "Elefanten" verstanden. Das ist die klassische Themafunktion.

\textbf{2. Kontrast}
\par\quad コーヒーは飲みますが、紅茶は飲みません。\\
\par\quad Kaffee trinke ich (schon), Tee (aber) nicht.\\

Hier hebt は zwei Elemente jeweils als Pole eines Gegensatzes heraus. Auch das ist Herausgreifung — nur mit einem impliziten Gegenstück statt einem neutralen Rahmen.

Beide Funktionen teilen denselben Kern: は macht etwas explizit zum Ausgangspunkt der Aussage und lässt dabei stets etwas Ungesagtes mitschwingen — beim Thema den Rahmen, beim Kontrast das ausgesparte Gegenstück.

\textbf{3. Warum が bei Fragewörtern steht, nicht は}
\par\quad 誰が来ますか。\\
\par\quad Wer kommt?\\

Ein Fragewort kann nicht herausgegriffen werden — es ist per Definition unbekannt, es gibt noch keinen Rahmen und keinen Kontrast, den man setzen könnte, bevor die Antwort feststeht. Deshalb steht hier immer が (die Identifikations-Partikel), nie は.

\section*{Eine ehrliche Grenze}

Bei Mengenangaben zeigt sich ein Fall, der sich nicht ganz sauber ins Bild einordnen lässt:

\par\quad 20人は集まりました。\\
\par\quad Es kamen mindestens 20 Leute (vielleicht mehr).\\

Hier markiert は eine Art Schwellenwert mit impliziter Offenheit nach oben — verwandt mit der Kontrast-Funktion, aber nicht restlos deckungsgleich mit ihr. Das ist keine Schwäche, die spezifisch unser Bild betrifft: Selbst in der Fachliteratur gibt es keine vollständig einheitliche Erklärung für alle Verwendungen von は. Das Herausgreifungs-Bild erklärt die weitaus meisten und wichtigsten Fälle zuverlässig — nur an diesem Rand bleibt eine kleine, offen eingestandene Unschärfe.


\chapter{の — Die Sach-Partikel}

\section*{Das Bild}

Stell dir の als ein unsichtbares Wort für "Ding" oder "Sache" vor, das sich an alles Mögliche anhängen lässt. の verbindet nicht nur zwei Nomen (wie ein Genitiv), es kann auch einen ganzen Satz oder ein Adjektiv in ein generisches "das X-Ding" verwandeln.

Diese beiden Funktionen wirken auf den ersten Blick verschieden, sind aber ein und dasselbe: の nimmt etwas — ein Nomen, einen Satz, eine Handlung — und macht daraus "die Sache von X" oder "das, was X betrifft".

\section*{Die wichtigsten Fälle}

\textbf{1. Genitiv-Verbindung}
\par\quad 私の本\\
\par\quad Mein Buch (wörtlich: das Buch der Sache "ich")\\

Das klassische の — verbindet zwei Nomen als Zugehörigkeit.

\textbf{2. Nominalisierung}
\par\quad 買ったのはこれです。\\
\par\quad Was ich gekauft habe, ist das hier. (wörtlich: "die Sache des Gekauft-Habens")\\

Hier verwandelt の einen ganzen Satz in ein "Ding", das dann wie ein Nomen weiterverwendet werden kann — als Subjekt, Objekt, oder womit auch immer der Satz weitermacht.

\textbf{3. Erklärender Satzabschluss}
\par\quad どうしたの？\\
\par\quad Was ist los? (wörtlich: "was ist die Sache, dass...")\\

の am Satzende deutet an: hier steckt ein Sachverhalt dahinter, der erklärt oder erfragt wird — dieselbe Nominalisierungs-Logik, nur am Satzende statt mitten im Satz.

\section*{Warum beide Funktionen zusammengehören}

Historisch gilt als wahrscheinlich, dass die Nominalisierungs-Funktion aus der Genitiv-Funktion durch eine Auslassung entstand: N1のもの ("das Ding, das zu N1 gehört") verlor im Lauf der Zeit vermutlich das Wort もの ("Ding") — übrig blieb N1の allein, das jetzt selbst die Bedeutung "das zugehörige Ding" trug. Beide Verwendungen wären demnach historisch dieselbe Konstruktion, nur an unterschiedlichen Stellen im Satz eingesetzt.

Das erklärt auch, warum の als Baustein in anderen Konstruktionen auftaucht: のに (obwohl) und ので (weil) bauen beide auf dem nominalisierten の auf — ein Satz wird zunächst zur "Sache" gemacht, an die dann に oder で andocken, mit ihren jeweils eigenen Bildern (Bezugspunkt bzw. Grenze/Scheidepunkt).

\section*{Ein praktischer Trick}

の lässt sich gedanklich fast immer durch もの (dingliche Sache) oder こと (Sachverhalt) ersetzen — das macht die Nominalisierungs-Funktion sofort greifbar:

\par\quad 赤いの = 赤いもの\\
\par\quad Das Rote (wörtlich: die rote Sache)\\

Das hilft auch bei der Frage, wann man の und wann こと als Nominalisierer wählt — es ist letztlich eine Wahl zwischen den beiden Nomen, für die の steht:

\par\quad 音楽を聞くのが好きです。\\
\par\quad Ich höre gern Musik. (の ≈ もの, konkreter/erlebnisnäher)\\

\par\quad 音楽を聞くことが好きです。\\
\par\quad Ich mag es, Musik zu hören. (こと, abstrakter/sachverhaltsbezogener)\\

Beide Sätze sind praktisch synonym, aber の betont eher das konkrete Erlebnis, こと eher den abstrakten Sachverhalt — passend dazu, dass の historisch für もの steht und こと für こと selbst.


\chapter{と — Die Gleichsetzungs-Partikel}

\section*{Das Bild}

Stell dir と als ein Gleichheitszeichen vor: zwei Dinge werden auf dieselbe Ebene gestellt, gleichrangig nebeneinander. と ist mehr als nur "mit" — es verbindet immer zwei Elemente, die in irgendeiner Weise als gleichwertig oder deckungsgleich gelten.

\section*{Die wichtigsten Fälle}

\textbf{1. Partner einer gemeinsamen Handlung}
\par\quad 友達とカラオケに行きます。\\
\par\quad Ich gehe mit einem Freund zum Karaoke.\\

Zwei gleichrangige Mitspieler derselben Handlung.

\textbf{2. Resultat einer Veränderung}
\par\quad 氷が水となります。\\
\par\quad Eis wird zu Wasser.\\

Hier wird der Endzustand mit dem Ausgangspunkt gleichgesetzt — nach der Veränderung \emph{ist} das eine das andere.

\textbf{3. Zitat}
\par\quad 「こんにちは」と言いました。\\
\par\quad Er sagte "Hallo".\\

Das Gesagte und der Wortlaut werden als exakt gleich gesetzt — と markiert die Äquivalenz zwischen Aussage und ihrem genauen Inhalt.

\textbf{4. Aufzählung}
\par\quad 新聞と雑誌\\
\par\quad Zeitung und Zeitschrift\\

Die Elemente stehen gleichrangig nebeneinander — die Reihenfolge lässt sich vertauschen (雑誌と新聞 funktioniert genauso gut), was die Gleichrangigkeit strukturell bestätigt.

\section*{Ein besonders enger Fall}

Bei bestimmten Verben wie 結婚する (heiraten) oder 争う (streiten) ist der と-Partner nicht optional wie bei bloßer Begleitung — er ist grammatisch notwendig, weil die Handlung wesensmäßig zwei gleichrangige Beteiligte braucht:

\par\quad 彼女と結婚します。\\
\par\quad Ich heirate sie.\\

Man kann hier nicht einfach mit と一緒に ("zusammen mit") ersetzen — die Handlung "Heiraten" existiert nur zwischen zwei gleichgestellten Partnern. と-Partner sind hier nicht bloß Begleiter, sondern strukturell gleichrangige Mit-Beteiligte der Handlung selbst.


\chapter{や — Die Beispiel-Partikel}

\section*{Das Bild}

Stell dir や als eine Handbewegung vor, die auf ein paar Dinge zeigt und sagt: "so etwas in der Art, und mehr davon." や zählt Beispiele auf, ohne zu behaupten, dass die Liste vollständig ist.

\section*{Der wichtigste Fall}

\par\quad 机の上に本や雑誌があります。\\
\par\quad Auf dem Tisch liegen Bücher, Zeitschriften und Ähnliches.\\

Hier klingt mit: es liegen wahrscheinlich noch mehr Dinge auf dem Tisch, aber Bücher und Zeitschriften sind repräsentative Beispiele.

\section*{Der Unterschied zu と}

In Kapitel と haben wir gesehen: と setzt zwei oder mehr Elemente vollständig und gleichrangig nebeneinander — die Liste ist damit abgeschlossen.

\par\quad 本と雑誌があります。\\
\par\quad Es sind ein Buch und eine Zeitschrift da. (nur diese beiden)\\

\par\quad 本や雑誌があります。\\
\par\quad Es sind Bücher, Zeitschriften und dergleichen da. (vermutlich noch mehr)\\

Derselbe Satzbau, aber und と und や setzen unterschiedliche Erwartungen: と schließt die Liste, や lässt sie offen.

\section*{Eine kleine Unschärfe}

In manchen Kontexten kann や auch verwendet werden, obwohl tatsächlich alle relevanten Elemente genannt sind — die Offenheit ist dann eher eine sprachliche Gewohnheit als eine strikte Aussage über Vollständigkeit. Das ist keine Ausnahme von der Grundidee, sondern zeigt nur: や \emph{impliziert} Offenheit, es \emph{garantiert} sie nicht in jedem einzelnen Fall.


\chapter{とか — Die Beispiel-Partikel mit Zurückhaltung}

\section*{Das Bild}

とか teilt den Grundgedanken mit や: Beispiele nennen, ohne Vollständigkeit zu behaupten. Aber とか geht weiter — es klingt nach einem beiläufigen, fast schulterzuckenden "so etwas in der Art, keine Ahnung, was noch alles". Es signalisiert nicht nur offene Aufzählung, sondern bewusste Unbestimmtheit.

\section*{Die wichtigsten Fälle}

\textbf{1. Aufzählung, informeller als や}
\par\quad 週末は映画とか見に行きたいな。\\
\par\quad Am Wochenende würde ich gern ins Kino oder so gehen.\\

Beispielhaft, locker, umgangssprachlich.

\textbf{2. Anschluss an mehr als Nomen}
\par\quad 忙しいとか疲れたとか、いろいろ理由をつける。\\
\par\quad Er findet allerlei Ausreden — beschäftigt, müde, was auch immer.\\

Anders als や, das nur an Nomen anschließt, funktioniert とか auch mit Adjektiven und Verben — es ist grammatisch flexibler.

\textbf{3. Bewusste Zurückhaltung}
\par\quad 仕事とか大変なんだ。\\
\par\quad Die Arbeit ist unter anderem stressig...\\

Hier klingt mehr mit: es gibt noch andere Dinge, über die der Sprecher nicht näher sprechen möchte. Interessanterweise taucht dieser Zurückhaltungs-Effekt manchmal sogar auf, wenn eigentlich schon alle relevanten Punkte genannt wurden (いいとか悪いとか, "gut oder schlecht oder so") — とか wird dann verwendet, um bewusst vage zu klingen, nicht weil tatsächlich mehr fehlt.

\section*{Der Unterschied zu や}

Beide Partikel teilen den Kern (offene Aufzählung), aber とか ist umgangssprachlicher — gegenüber Vorgesetzten oder in formeller Sprache eher zu vermeiden — und trägt diesen zusätzlichen Ton der beiläufigen Zurückhaltung, den や nicht hat.


\chapter{も — Die Akkumulations-Partikel}

\section*{Das Bild}

Stell dir も als ein Plus-Zeichen vor: "und das hier auch." も hebt ein Element heraus und reiht es zu einer Menge anderer, für die dasselbe bereits gilt. Wo は einen Kontrast setzt (dies schon, das nicht) und が identifiziert (genau das, und nur das), addiert も (auch das noch, zusätzlich zum Bekannten).

\section*{Die wichtigsten Fälle}

\textbf{1. Einfache Addition}
\par\quad 太郎も学生です。\\
\par\quad Auch Taro ist Student.\\

Taro wird zu einer Gruppe (anderer Studenten) hinzugefügt, für die dasselbe Prädikat gilt.

\textbf{2. Universalquantor mit Verneinung}
\par\quad 誰も来ませんでした。\\
\par\quad Niemand kam.\\

Hier akkumuliert も über \emph{alle} möglichen Personen hinweg — "nicht dieser, nicht jener, nicht..." bis zur vollständigen Verneinung. Dieselbe Addition, nur konsequent bis zum Ende durchgezogen.

\textbf{3. Mengenbetonung}
\par\quad ビールを10本も飲みました。\\
\par\quad Er hat sage und schreibe 10 Bier getrunken.\\

も signalisiert hier: diese Menge übersteigt die Erwartung — auch das ist Akkumulation, nur gegen eine implizite Erwartungsschwelle statt gegen eine Personengruppe.

\textbf{4. Extrembeispiel}
\par\quad お茶でも飲みますか。\\
\par\quad Möchten Sie vielleicht Tee (oder etwas Ähnliches) trinken?\\

Hier wird ein naheliegendes Beispiel genannt, mit dem Unterton "und alles andere in dieser Richtung würde auch gehen" — die Akkumulation bleibt implizit, deutet aber auf eine offene Menge von Alternativen.

\section*{Ein kombinierter Fall: でも}

でも setzt sich zusammen aus で (der Grenz-Partikel, die einen Bereich markiert) und も (Akkumulation). Das erklärt seine Verwendungen direkt:

\par\quad 誰でも来られます。\\
\par\quad Jeder kann kommen (egal wer).\\

Hier wird nicht nur über Personen akkumuliert (誰も), sondern über \emph{jeden denkbaren Fall/Bereich} — で liefert den Bereich, も die Akkumulation darüber. でも braucht kein eigenes Bild, es ergibt sich vollständig aus seinen beiden Bausteinen.


\chapter{から — Die Ausgangspunkt-Partikel}

\section*{Das Bild}

Stell dir から als einen Startpunkt auf einer Linie vor — der Ort, die Zeit, der Zustand oder der Grund, von dem aus sich etwas entwickelt. から markiert immer, wo etwas seinen Anfang nimmt.

\section*{Die wichtigsten Fälle}

\textbf{1. Räumlicher Ausgangspunkt}
\par\quad 駅から電車で向かいます。\\
\par\quad Ich fahre vom Bahnhof aus mit dem Zug.\\

\textbf{2. Ausgangszustand einer Veränderung}
\par\quad 水から氷になりました。\\
\par\quad Aus Wasser wurde Eis.\\

Hier steht から für den Zustand \emph{vor} der Veränderung — das Gegenstück zu と, das (aus Kapitel と) den \emph{Endzustand} mit dem Ergebnis gleichsetzt (氷が水となる). から blickt zurück auf den Ursprung, と markiert das Resultat.

\textbf{3. Grund/Ursache}
\par\quad 忙しいから、行けません。\\
\par\quad Weil ich beschäftigt bin, kann ich nicht kommen.\\

Der Grund wird als Ausgangspunkt der Schlussfolgerung gedacht — von dieser Tatsache aus ergibt sich das Resultat.

\textbf{4. Aufeinanderfolgende Handlungen}
\par\quad 食べてから、寝ます。\\
\par\quad Nachdem ich gegessen habe, schlafe ich.\\

Die erste, abgeschlossene Handlung wird zum Ausgangspunkt für die zweite.

\section*{Der Unterschied zu に bei Empfängern}

Sowohl に als auch から können bei Verben wie もらう ("bekommen") den Geber markieren:

\par\quad 友達に本をもらいました。\\
\par\quad 友達から本をもらいました。\\
\par\quad Ich habe von einem Freund ein Buch bekommen.\\

Der Unterschied: に betont den Freund als \textbf{Interaktionspartner} (Bezugspunkt einer Beziehung), から betont ihn als \textbf{reinen Ursprung} — deshalb wird bei unpersönlichen Quellen wie Institutionen fast immer から verwendet (組織から奨学金をもらう, "ein Stipendium vom Ministerium bekommen"), nicht に: eine Organisation ist kein Interaktionspartner im relationalen Sinn, sondern schlicht eine Quelle.


\chapter{まで — Die Endpunkt-Partikel}

\section*{Das Bild}

まで markiert den Endpunkt einer durchgehenden Erstreckung — den Punkt, bis zu dem etwas andauert, wobei dieser Punkt selbst noch dazugehört. Wo から den Anfang setzt, setzt まで das Ende, und beide zusammen spannen eine Strecke auf.

\section*{Die wichtigsten Fälle}

\textbf{1. Räumlicher/zeitlicher Endpunkt}
\par\quad 5ページから10ページまで読みました。\\
\par\quad Ich habe von Seite 5 bis Seite 10 gelesen.\\

Seite 10 gehört noch mit dazu — まで schließt den Endpunkt ein.

\textbf{2. Inklusiver Zeitrahmen}
\par\quad 15日まで待ってください。\\
\par\quad Bitte warten Sie bis zum 15.\\

Der 15. selbst ist noch eingeschlossen.

\section*{Der feine Unterschied zu までに}

\par\quad 5時まで待ちます。\\
\par\quad Ich warte bis 5 Uhr (durchgehend, die ganze Zeit).\\

\par\quad 5時までに終わります。\\
\par\quad Es ist bis spätestens 5 Uhr fertig (ein Termin, kein durchgehender Zeitraum).\\

まで allein beschreibt eine Fortdauer — die Handlung läuft ununterbrochen bis zu diesem Punkt. Kommt に dazu (までに), wird daraus ein Termin: nicht mehr die Erstreckung selbst, sondern ein einzelner Punkt, den man bis dahin erreichen muss. Das に bringt hier genau die Punkt-Logik ein, die wir schon aus dem に-Kapitel kennen.

\section*{Die Klammer mit から}

から und まで gehören zusammen: から markiert, wo die Strecke beginnt, まで, wo sie endet.

\par\quad 東京から大阪まで新幹線で行きます。\\
\par\quad Ich fahre mit dem Shinkansen von Tokyo bis Osaka.\\

Zusammen bilden sie die vollständige Klammer einer durchmessenen Erstreckung — passend zu を, das wir aus dem Kapitel zu dieser Partikel kennen: を markiert das, was vollständig durchmessen wird, から...まで markiert, wo diese Durchmessung beginnt und endet.


\chapter{だけ — Die Exklusivitäts-Partikel}

\section*{Das Bild}

だけ zieht eine scharfe, aber wertungsfreie Grenze um eine Menge oder Auswahl: "genau das, nicht mehr, nicht weniger." Kein Bedauern, kein Erstaunen — nur eine schlichte Feststellung der Grenze.

\section*{Der wichtigste Fall}

\par\quad 一人だけ来ました。\\
\par\quad Nur eine Person kam.\\

Neutral konstatiert: genau eine, nicht mehr.

\section*{Der Unterschied zu しか}

Japanisch hat mehrere Wörter für "nur", die sich in ihrem Unterton unterscheiden. だけ ist der neutrale unter ihnen:

\par\quad 一人だけ来ました。\\
\par\quad Nur eine Person kam. (schlichte Feststellung)\\

\par\quad 一人しか来ませんでした。\\
\par\quad Nur eine einzige Person kam (und das ist zu wenig).\\

しか verlangt zwingend eine Verneinung im Satz und trägt immer einen Unterton der Unzulänglichkeit — "das ist zu wenig, schade". だけ dagegen bleibt neutral: eine reine Mengenangabe ohne emotionale Färbung.

\bigskip\hrule\bigskip

\emph{Für Interessierte:} だけ geht auf das alte Nomen たけ (丈) zurück, das "Maß" oder "Ausdehnung" bedeutete — etwa in 背丈 ("Körpergröße", wörtlich "Rückenmaß"). "Bis zum Maß von X" passt genau zum heutigen Bild: だけ zieht die Grenze exakt am Maß dessen, was gilt.


\chapter{より — Die Vergleichs-Ausgangspunkt-Partikel}

\section*{Das Bild}

より sieht wie eine eigenständige "als"-Partikel aus, ist aber eigentlich eine Spezialisierung von から: statt eines räumlichen oder zeitlichen Ausgangspunkts markiert より den \textbf{Ausgangspunkt einer Vergleichsachse}.

\section*{Beispiel}

\par\quad 東京は大阪より大きいです。\\
\par\quad Tokyo ist größer als Osaka.\\

Wörtlich gedacht: "Tokyo ist groß, gemessen ausgehend von Osaka." Osaka ist der Referenzpunkt, von dem aus der Größenunterschied gemessen wird — genau wie から einen Ausgangspunkt für eine Bewegung markiert, markiert より einen Ausgangspunkt für einen Vergleich.

\section*{Warum das keine Metapher ist}

Anders als bei den meisten Partikelbildern ist der Zusammenhang zwischen より und から hier nicht nur eine gedankliche Brücke, sondern historisch direkt belegt: より war im klassischen Japanisch tatsächlich die primäre Ausgangspunkt-Partikel, aus der sich die heutige Vergleichsfunktion erst später als Spezialfall entwickelte. Die alte Ausgangspunkt-Bedeutung lebt in formeller/schriftlicher Sprache bis heute weiter:

\par\quad 5ページより先を読んでください。\\
\par\quad Bitte lesen Sie ab Seite 5 weiter.\\

Hier hat より exakt dieselbe Funktion wie から — nur in gehobenerem Register.

\section*{Eine Sonderkonstruktion}

\par\quad これより他に方法がない。\\
\par\quad Es gibt keine andere Möglichkeit als diese.\\

Auch hier ist より der Ausgangspunkt: "ausgehend von dieser Möglichkeit gibt es keine andere" — derselbe Gedanke, nur mit Verneinung kombiniert zu einer Ausschluss-Konstruktion.


\chapter{し — Die Akkumulierende Grund-Partikel}

\section*{Das Bild}

し funktioniert wie も, nur auf Satzebene: während も ein Nomen zu einer Menge anderer addiert ("auch das noch"), addiert し einen Sachverhalt oder Grund zu weiteren, die mitgelten — "und außerdem noch das."

\section*{Die wichtigsten Fälle}

\textbf{1. Mehrere Gründe aufgezählt}
\par\quad 物価も高いし、電車も人が多いし、それに暑い。\\
\par\quad Die Preise sind hoch, die Bahnen sind überfüllt, und dazu ist es heiß.\\

Jeder Grund wird mit し angehängt — die Aufzählung bleibt offen, es könnten noch mehr Gründe folgen.

\textbf{2. Ein einzelner Grund, mit impliziten weiteren}
\par\quad 疲れたし、もう寝ます。\\
\par\quad Ich bin müde, also gehe ich jetzt schlafen.\\

Auch wenn nur ein Grund genannt wird, schwingt mit し oft mit: es gibt wahrscheinlich noch mehr Gründe, die einfach nicht ausgesprochen werden. Das unterscheidet し von から, das einen Grund ohne diesen Anspruch auf weitere Gründe nennt — から ist die reine, einzelne Begründung, し deutet eine ganze Sammlung an.

\section*{Warum し zurückhaltender wirkt als から}

Weil し immer eine offene Menge an mitgeltenden Gründen andeutet, wirkt eine Begründung mit し oft weniger kategorisch oder endgültig als eine mit から — es klingt eher nach "unter anderem, weil..." als nach "genau deshalb, weil...".


\chapter{けれど — Die Konzessions-Partikel}

\section*{Das Bild}

けれど verbindet zwei Satzteile, bei denen der zweite gegen die Erwartung verstößt, die der erste weckt: "Obwohl das eine gilt, gilt trotzdem das andere." Es ist die gesprochene, alltägliche Variante von "obwohl" oder "aber".

\section*{Der wichtigste Fall}

\par\quad 高いけれど、買いました。\\
\par\quad Es war teuer, aber ich habe es trotzdem gekauft.\\

Der erste Teil (teuer) lässt eine bestimmte Konsequenz erwarten (nicht kaufen) — der zweite Teil widerspricht dieser Erwartung.

\section*{Verwandtschaft zu のに}

Diese Funktion kennen wir bereits aus dem Kapitel の: die Konstruktion のに funktioniert nach demselben Muster — ein Satz weckt eine Erwartung, die vom folgenden Satz nicht erfüllt wird. けれど ist gewissermaßen die eigenständige Konjunktionsvariante desselben konzessiven Gedankens, während のに direkt an die Nominalisierungsfunktion von の anknüpft.

\section*{Register}

けれど kommt in verschiedenen Kürzungsstufen vor — けれども (am formellsten), けれど, けど (am beiläufigsten) — je nachdem, wie locker die Sprechsituation ist. Die Bedeutung bleibt in allen Formen dieselbe.


\chapter{ながら — Die Koexistenz-Partikel}

\section*{Das Bild}

ながら sagt: zwei Dinge bestehen gleichzeitig, ohne sich gegenseitig aufzuheben. Am bekanntesten ist die Verwendung für gleichzeitige Handlungen, aber dieselbe Grundidee trägt noch drei weitere Fälle.

\section*{Die wichtigsten Fälle}

\textbf{1. Gleichzeitige Handlungen}
\par\quad 音楽を聴きながら勉強します。\\
\par\quad Ich lerne, während ich Musik höre.\\

Zwei Handlungen laufen nebeneinander her.

\textbf{2. Konzession — zwei koexistierende, unvereinbar wirkende Zustände}
\par\quad 貧しいながらも幸せな家庭です。\\
\par\quad Eine Familie, die arm und doch glücklich ist.\\

Hier bestehen "arm" und "glücklich" gleichzeitig, obwohl man das eine eher nicht mit dem anderen erwarten würde. Die "obwohl"-Bedeutung entsteht ganz natürlich aus der Koexistenz: wenn zwei gleichzeitig bestehende Dinge in Spannung zueinander stehen, liest sich das automatisch als Widerspruch.

\textbf{3. Zustandsfortdauer}
\par\quad 昔ながらのたたずまい。\\
\par\quad Ein Erscheinungsbild wie schon immer (unverändert seit früher).\\

Der alte Zustand koexistiert bis in die Gegenwart hinein — nichts hat sich verändert.

\textbf{4. Totalität}
\par\quad 三度ながら失敗しました。\\
\par\quad Alle drei Male bin ich gescheitert.\\

Mehrere Vorkommnisse werden als eine zusammengehaltene Einheit gedacht — auch das eine Form von Koexistenz: die drei Male "bestehen zusammen" als ein Ganzes.

\section*{Eine Faustregel}

Bei Bewegungsverben (音楽を聴きながら) liest sich ながら meist als schlichte Gleichzeitigkeit. Bei Adjektiven oder Zustandswörtern (貧しいながら) liest es sich eher konzessiv. Das ist keine strikte Regel, aber eine brauchbare erste Orientierung.


\chapter{か — Die Ungewissheits-Partikel}

\section*{Das Bild}

か markiert, dass eine Information nicht feststeht — sie liegt außerhalb der Kontrolle oder Gewissheit des Sprechers. Die bekannteste Verwendung ist die Frage, aber das ist nur der sichtbarste Fall einer viel allgemeineren Grundidee.

\section*{Die wichtigsten Fälle}

\textbf{1. Frage}
\par\quad あなたも行きますか。\\
\par\quad Gehen Sie auch?\\

Die Information liegt beim Hörer — der Sprecher weiß es nicht.

\textbf{2. Unbestimmtheit bei Fragewörtern}
\par\quad 誰かが来ました。\\
\par\quad Irgendjemand ist gekommen.\\

Hier bleibt bewusst offen, wer genau gemeint ist — か hängt sich an das Fragewort und macht daraus eine unbestimmte Angabe statt einer echten Frage.

\textbf{3. Alternative}
\par\quad コーヒーか紅茶か選んでください。\\
\par\quad Bitte wählen Sie Kaffee oder Tee.\\

Welche der beiden Optionen zutrifft, bleibt offen — か markiert genau diese Unentschiedenheit.

\textbf{4. Selbstfrage}
\par\quad どうしようか。\\
\par\quad Was soll ich nur tun?\\

Interessant: hier behandelt der Sprecher sogar die eigene, noch zu treffende Entscheidung als ungewiss — か zeigt, dass selbst der Sprecher selbst die Antwort noch nicht kennt.

\textbf{5. Überraschung über neue Information}
\par\quad だれかと思ったら、君だったのか。\\
\par\quad Ach, du warst es!\\

Eine gerade erst erfahrene Tatsache, die noch nicht vollständig "verarbeitet" oder integriert ist — auch das ist eine Form von Ungewissheit, nur bezogen auf den Moment der Überraschung.

\section*{Ein größeres Muster}

か fügt sich in eine Reihe von Partikeln ein, die wir bereits kennen und die alle mit dem Grad der Gewissheit einer Aussage zu tun haben: と (Kapitel と) legt eine vollständige, feste Aufzählung fest, や (Kapitel や) lässt bewusst offen, ob noch mehr dazugehört, とか (Kapitel とか) fügt dem noch eine bewusste Zurückhaltung hinzu. か geht am weitesten: hier steht nicht nur die Vollständigkeit einer Aufzählung offen, sondern die Information selbst.


\chapter{ね — Die Partikel geteilter Gewissheit}

\section*{Das Bild}

ね sagt: "Das wissen wir doch beide, oder?" Es markiert eine Information als gemeinsamen Besitz von Sprecher und Hörer und sucht Bestätigung oder Zustimmung dafür. ね rückt die Gesprächspartner näher zusammen — es lädt den Hörer ein, mitzuschwingen.

\section*{Der wichtigste Fall}

\par\quad いい天気ですね。\\
\par\quad Schönes Wetter, nicht wahr?\\

Der Sprecher geht davon aus, dass der Hörer dasselbe sieht und fühlt, und sucht eine Bestätigung dieser geteilten Wahrnehmung.

\par\quad もうすぐ夏休みですね。\\
\par\quad Bald sind Sommerferien, nicht wahr?\\

Auch diese Tatsache kennen beide Gesprächspartner ohnehin — sie wird gemeinsam bestätigt, nicht um Information zu liefern, sondern um Nähe und Übereinstimmung herzustellen.


\chapter{よ — Die Partikel exklusiven Sprecherwissens}

\section*{Das Bild}

よ sagt: "Das weißt du noch nicht, aber ich sage es dir jetzt." Anders als ね, das ein geteiltes Wissen bestätigt, teilt よ eine Information mit, die (nach Annahme des Sprechers) allein bei ihm liegt und dem Hörer neu ist. よ trägt dabei oft auch Nachdruck — der Sprecher besteht auf seiner Sichtweise.

\section*{Die wichtigsten Fälle}

\textbf{1. Neue Information mitteilen}
\par\quad わたしは行きませんよ。\\
\par\quad Ich gehe nicht, das sage ich dir!\\

Der Sprecher teilt eine Tatsache mit, die der Hörer noch nicht kennt.

\textbf{2. Aufmerksamkeit lenken, Nachdruck verleihen}
\par\quad 危ないよ！\\
\par\quad Vorsicht!\\

Hier verstärkt よ die Dringlichkeit — der Sprecher will sicherstellen, dass die Warnung ankommt.

\section*{Der Unterschied zu ね}

In Kapitel ね haben wir gesehen: ね setzt geteiltes Wissen voraus und sucht Bestätigung dafür. よ dreht das um — es teilt etwas mit, das der Sprecher exklusiv zu wissen glaubt, und der Hörer noch nicht kennt.

\par\quad 明日、雨ですよ。\\
\par\quad Morgen regnet es, das wusstest du noch nicht.\\

\par\quad 明日、雨ですね。\\
\par\quad Morgen regnet es, nicht wahr? (das wissen wir doch schon beide)\\

Beide Partikel beantworten dieselbe stille Frage — "wo liegt dieses Wissen, bei mir allein oder bei uns beiden?" — nur mit entgegengesetzter Antwort.


\chapter{かい — Die familiäre Frage-Variante}

\section*{Das Bild}

かい setzt sich zusammen aus か (der Ungewissheits-Partikel, die wir bereits kennen) und い — zwei Fragebausteine, kombiniert zu einer besonders vertrauten, umgangssprachlichen Fragevariante. Eine eigenständige Bedeutung über か hinaus hat かい nicht; es ist im Kern dieselbe Ungewissheits-Markierung, nur mit einem zusätzlichen Ton von Vertrautheit oder Beiläufigkeit.

\section*{Der wichtigste Fall}

\par\quad 元気かい？\\
\par\quad Alles gut bei dir?\\

Dieselbe Frage wie 元気ですか, aber deutlich lockerer, familiärer im Ton — typisch unter engen Freunden oder in sehr informeller Sprache.

\section*{Eine Klarstellung zum Sprachgebrauch}

かい wird oft als reine Männersprache beschrieben. Das trifft die heutige Wahrnehmung durchaus — in der Popkultur (etwa in Yakuza-Filmen) klingt かい stark maskulin. Historisch war かい aber nicht ausschließlich auf männliche Sprecher beschränkt; die strikte Zuordnung ist eher ein modernes Klischee als eine feste Sprachregel. In der Praxis heute gilt trotzdem: かい wirkt in den meisten Kontexten männlich-umgangssprachlich und sollte mit Bedacht verwendet werden.


\chapter{だい — Die familiäre Wortfrage-Variante}

\section*{Das Bild}

だい ist das Gegenstück zu かい aus dem vorigen Kapitel — dieselbe Idee einer weichen, vertrauten Fragevariante, nur für eine andere Art von Frage. Während かい bei einfachen Ja/Nein-Fragen steht, gehört だい zu Fragen mit einem Fragewort (wer, was, wo, wann...).

\section*{Der wichtigste Fall}

\par\quad どこに行くんだい？\\
\par\quad Wo gehst du denn hin?\\

\par\quad それ、何だい？\\
\par\quad Was ist das denn?\\

Dieselbe Frage ließe sich auch neutral mit か stellen (どこに行くんですか) — だい macht sie spürbar lockerer und vertrauter, typisch für ältere, familiäre Sprechweise.

\section*{Die Arbeitsteilung mit かい}

Beide Formen — かい für Ja/Nein-Fragen, だい für Fragewort-Fragen — funktionieren nach demselben Prinzip: ein zusätzlicher Vokal macht eine direkte Frage weicher und vertrauter. かい hängt sich an か, だい an だ.

\section*{Eine Anmerkung zum Sprachgebrauch}

だい wirkt in der Praxis überwiegend männlich-umgangssprachlich, mit der bekannten Ausnahme, dass es traditionell auch bei älteren Frauen vorkommt. Wie bei かい gilt: mit Bedacht verwenden, da die Form einen deutlichen Registerton trägt.


\chapter{Die Verb-Stammformen — der Baukasten hinter den Formen}

\section*{Ein anderes Kapitel als bisher}

Bislang ging es immer um eine Partikel — eine Wahlmöglichkeit, bei der ein Sprecher sich für ein bestimmtes Bild entscheidet. Die Stammformen sind etwas anderes: es sind die Bausteine, aus denen sich fast alle Verbformen zusammensetzen — Verneinung, Höflichkeitsform, Befehlsform, Bedingungssatz und vieles mehr. Ein Verb wird nicht in einer Stammform "verwendet", sondern durch sie hindurch zu einer anderen Form weitergebaut.

\section*{Der praktische Kern: die Endvokale}

Bei Verben mit Stamm auf Konsonant ("u-Verben", z.B. 書く) erkennt man die Stammform direkt am Endvokal. In der japanischen Vokalreihenfolge (a–i–u–e–o):

| Vokal | Stammform | Beispiel (書く) | Führt zu |
|---|---|---|---|
| -a | 未然形 | 書か | ない-Form, せる (Kausativ), れる (Passiv) |
| -i | 連用形 | 書き | ます-Form, たい-Form |
| -u | 終止形 | 書く | Grundform (Wörterbuchform) |
| -e | 仮定形/命令形 | 書け | ば-Form (Bedingung), Befehlsform |
| -o | 未然形2/推量形 | 書こ | よう-Form (Vermutung/Wille, +う) |

音便形 hat keinen eigenen Endvokal — sie entsteht als Lautveränderung der 連用形 vor て/た (書き+て → 書いて).

\section*{Andere Verbgruppen}

Bei Verben, deren Stamm auf einen Konsonanten endet (z.B. 書く, 話す, 読む — oft "u-Verben" genannt, weil die Grundform auf -u endet), sieht man den Vokalwechsel deutlich, weil der Stamm selbst den Vokal trägt. Bei den anderen Verbgruppen ist das anders:

\textbf{Verben mit Stamm auf -eru/-iru (z.B. 食べる, 見る)} — der Stamm bleibt unverändert, nur die Endung ändert sich:

| Stammform | 食べる |
|---|---|
| 未然形 | 食べ (+ない, +させる, +られる) |
| 連用形 | 食べ (+ます, +たい) |
| 終止形 | 食べる |
| 仮定形/命令形 | 食べれ (+ば) / 食べろ |
| 未然形2/推量形 | 食べよ (+う) |

\textbf{Die zwei unregelmäßigen Verben} — する und 来る folgen keinem der beiden Muster durchgehend und müssen einzeln gelernt werden:

| Stammform | する | 来る |
|---|---|---|
| 未然形 | し (+ない), さ (+せる), さ (+れる) | 来 (こ) (+ない) |
| 連用形 | し (+ます, +たい) | 来 (き) (+ます, +たい) |
| 終止形 | する | 来る (くる) |
| 仮定形/命令形 | すれ (+ば) / しろ | 来れ (くれ) (+ば) / 来い (こい) |
| 未然形2/推量形 | し (+よう) | 来 (こ) (+よう) |

Sieben traditionelle Stammformen (eigentlich sechs plus eine spätere Auftrennung) bilden dieses Fundament. Bei manchen davon passen alle daraus gebildeten Formen zum selben Grundgedanken, bei anderen nur ein Teil — und gerade das macht sie interessant.

\section*{Die Bausteine im Überblick}

\textbf{連用形 (Verbindungsform)} — funktioniert wie の: sie macht aus dem Verb ein "Ding", das an anderes andocken kann (ますForm, たいForm, oder Verbindungen mit anderen Wörtern). Diese Nominalisierungs-Idee trägt konsequent durch alle Formen, die aus diesem Baustein entstehen.

\textbf{終止形 (Schlussform)} — die Grundform, mit der ein Satz vollständig abgeschlossen wird.

\textbf{未然形 (Noch-nicht-Form)} — hier wird es interessant: Die ないForm (Verneinung — etwas ist "noch nicht" der Fall) und die Vermutungsform passen zu diesem Gedanken. Kausativ (せる) und Passiv (れる), die ebenfalls aus diesem Baustein gebildet werden, tun das nicht — sie folgen ihm nur aus lautlichen Gründen (eine Bindevokal-Notwendigkeit), nicht aus einem gemeinsamen "noch nicht"-Gedanken.

\textbf{仮定形 (Hypothetische Form)} — ähnlich gemischt: Der Bedingungssatz mit ば trägt die "wenn"-Idee, andere daraus gebildete Formen tun das nicht in gleicher Weise.

\textbf{志向形 — zwei Formen, fälschlich als eine behandelt.} Was traditionell als eine Stammform galt, sind eigentlich zwei unabhängige Bausteine: die Befehlsform (命令形) und die Vermutungs-/Willensform (推量形, Basis für die ようForm). Sie haben nicht einmal denselben Stammvokal.

\textbf{音便形 (Lautwandel-Form)} — beschreibt einen rein lautlichen Vorgang: て-Form und た-Form entstehen aus demselben Aussprache-Erleichterungs-Mechanismus, ohne eigene Bedeutung.

\section*{Was das für die folgenden Kapitel bedeutet}

In den nächsten Kapiteln geht es um die konkreten Formen, die aus diesen Bausteinen entstehen — ますForm, たいForm, ないForm, Kausativ, Passiv und so weiter. Dort zählt jeweils die eigene Bedeutung der Form selbst, nicht, aus welchem Baustein sie technisch gebildet wird. Dieses Kapitel dient nur als Hintergrundwissen, warum bestimmte Formen morphologisch zusammenhängen, obwohl sie inhaltlich nichts miteinander zu tun haben.


\chapter{ないForm — Die Verneinung}

\section*{Das Bild}

ない markiert Abwesenheit — dass etwas nicht der Fall ist, nicht existiert oder nicht geschieht. Egal ob an ein Verb oder direkt als eigenständiges Wort verwendet, der Kerngedanke bleibt derselbe: etwas fehlt.

\section*{Die wichtigsten Fälle}

\textbf{1. Verbverneinung}
\par\quad 食べない。\\
\par\quad Ich esse nicht.\\

Die Handlung ist abwesend — sie findet nicht statt.

\textbf{2. Als Verneinung von ある: "nicht existieren"}
\par\quad お金がない。\\
\par\quad Ich habe kein Geld. (wörtlich: Geld existiert nicht.)\\

Wo ある ("existieren, vorhanden sein") die Existenz eines Dings behauptet, verneint ない sie — dieselbe Abwesenheits-Idee wie bei der Verbverneinung, nur auf ein Ding statt auf eine Handlung bezogen.

\textbf{3. Ratschlag}
\par\quad 早く寝た方がいいよ。\\
\par\quad Du solltest lieber früh schlafen gehen.\\

Die Verneinungsform steckt hier in der Vergleichskonstruktion (方がいい, "besser wäre..."), mit der man implizit auf die Abwesenheit der schlechteren Alternative hinweist.

\textbf{4. Im Begriff sein, etwas nicht zu tun}
\par\quad 食べないところだ。\\
\par\quad Ich bin dabei, nicht zu essen (ich esse gerade nicht).\\

Auch in dieser Konstruktion wird die Abwesenheit der Handlung zum zentralen Punkt der Aussage.

\section*{Eine kleine Besonderheit}

Die Verneinung von Verben (読まない, "nicht lesen") und die Verneinung von ある (お金がない, "kein Geld existiert") haben ursprünglich unterschiedliche Wurzeln — sie sind historisch nicht dasselbe Wort. Erst im Lauf der Zeit haben sie sich derselben Konjugationsweise angeglichen. Trotzdem tragen beide bis heute denselben Grundgedanken: etwas ist nicht da.


\chapter{たいForm — Der subjektive Wunsch}

\section*{Das Bild}

たい drückt einen Wunsch aus, den man von innen heraus kennt — den eigenen. Es ist die Form, mit der man sagt, was man selbst will, nicht was man bei anderen beobachtet.

\section*{Der wichtigste Fall}

\par\quad 日本に行きたいです。\\
\par\quad Ich möchte nach Japan fahren.\\

Ein Wunsch aus der eigenen, unmittelbaren Perspektive.

\section*{Die Grenze: nur für sich selbst (und im Gespräch für den Zuhörer)}

たい wird normalerweise nur für die eigene Person verwendet, oder in Fragen für den Gesprächspartner:

\par\quad 何が食べたいですか。\\
\par\quad Was möchtest du essen?\\

Für eine dritte Person passt たい dagegen nicht ohne Weiteres — man kann von außen nicht direkt in den Wunsch eines anderen hineinsehen, sondern nur sein Verhalten beobachten:

\par\quad 彼は日本に行きたがっています。\\
\par\quad Er scheint nach Japan fahren zu wollen. (beobachtet, nicht direkt gewusst)\\

Hier braucht es die Form たがる — sie beschreibt einen Wunsch, wie er von außen sichtbar wird, nicht wie er von innen erlebt wird. Dieser Unterschied — was der Sprecher direkt weiß, gegenüber dem, was er nur von außen beobachten kann — ist ein wiederkehrendes Prinzip im Japanischen: dieselbe Logik begegnet uns bei Adjektiven, die Gefühle beschreiben (z.B. 痛い, "schmerzhaft"), die für eine dritte Person ebenfalls die がる-Form annehmen (痛がる, "Schmerz zeigen").


\chapter{ますForm — Die Höflichkeitsform}

\section*{Das Bild}

ます macht aus einer neutralen Aussage eine höfliche — die Form, mit der man einem Gegenüber Respekt und Distanz entgegenbringt, ohne dabei kalt oder unpersönlich zu wirken. Es ist die Standard-Höflichkeitsebene im Japanischen, die man in den allermeisten Alltagssituationen mit Fremden, Kollegen oder Respektspersonen verwendet.

\section*{Der wichtigste Fall}

\par\quad 毎日、日本語を勉強します。\\
\par\quad Ich lerne jeden Tag Japanisch.\\

Dieselbe Aussage ohne ます (勉強する) wäre inhaltlich identisch, klingt aber direkt und familiär — passend unter Freunden, aber nicht im formellen Kontext.

\section*{Warum diese Form Höflichkeit ausdrückt}

ます geht vermutlich auf ein altes, bescheidenes Verb zurück, mit dem man ursprünglich signalisierte, dass man sich selbst gegenüber dem Adressaten einer Handlung zurücknimmt — die genaue Herkunft ist in der Forschung nicht einheitlich geklärt. Nach dieser Lesart entwickelte sich aus der ursprünglichen Bescheidenheit gegenüber der Person, auf die sich die Handlung bezog, mit der Zeit die heutige, allgemeinere Höflichkeit gegenüber dem Gesprächspartner selbst — unabhängig davon, worum es im Satz eigentlich geht.


\chapter{よう — Zwei Formen mit demselben Klang}

\section*{Eine Warnung vorweg}

よう sieht auf den ersten Blick wie eine einzige Form aus. Tatsächlich verbergen sich dahinter zwei völlig unabhängige Bausteine, die zufällig gleich klingen — ähnlich wie im Deutschen "Kiefer" (der Knochen im Gesicht) und "Kiefer" (der Nadelbaum), die trotz gleicher Schreibung auf ganz unterschiedliche Wurzeln zurückgehen und nichts miteinander zu tun haben. Wer beide für dieselbe Sache hält, gerät schnell durcheinander. Deshalb werden sie hier bewusst getrennt behandelt.

\section*{Die erste よう: Wille und Vermutung}

Diese Form entsteht direkt aus dem Verb selbst (行こう, 食べよう) und drückt Willensäußerung oder Vermutung aus:

\par\quad 一緒に行こう。\\
\par\quad Lass uns zusammen gehen!\\

\par\quad もうすぐ雨が降ろう。\\
\par\quad Es wird wohl bald regnen.\\

よう steht dabei nicht zwingend am Satzende — es kann auch mitten im Satz auftauchen, etwa vor と:

\par\quad 早く行こうと思っています。\\
\par\quad Ich denke daran, bald zu gehen.\\

\section*{Die zweite よう: Vergleich, Wahrscheinlichkeit, Zweck}

Diese Form geht auf das Nomen 様 ("Weise, Art") zurück und tritt meist als ようだ oder ように auf:

\par\quad 雨が降るようだ。\\
\par\quad Es sieht aus, als würde es regnen.\\

\par\quad 猫のように寝ています。\\
\par\quad Er schläft wie eine Katze.\\

\par\quad 忘れないように書いておきます。\\
\par\quad Ich schreibe es auf, damit ich es nicht vergesse.\\

Hier geht es nicht um Wille, sondern um Ähnlichkeit, Anschein oder Zweck — inhaltlich hat das nichts mit der ersten よう zu tun, auch wenn der Klang identisch ist.

\section*{Wo beide aufeinandertreffen}

Ein Fall macht die Verwechslungsgefahr besonders deutlich:

\par\quad どうしようもない。\\
\par\quad Da kann man nichts machen.\\

Hier klingt よう wie das Willens-よう (しよう, "ich werde tun"), gehört aber eigentlich zur zweiten Familie — 様 im Sinne von "Weise, Methode" ("es gibt keine Weise/Methode, es zu tun"). Genau solche Fälle zeigen, warum es sich lohnt, die beiden よう von Anfang an als zwei getrennte Dinge zu begreifen, auch wenn sie gleich klingen.


\chapter{そうForm — Eindruck und Hörensagen}

\section*{Das Bild}

そう markiert, dass eine Aussage nicht auf direktem eigenem Wissen beruht, sondern auf einem äußeren Eindruck oder auf dem, was man gehört hat. Es gibt zwei Verwendungen, die sich äußerlich sehr ähnlich sehen, aber an unterschiedlichen Stellen im Satz andocken — und genau daran erkennt man, welche gemeint ist.

\section*{Eindruck (直接andockend an den Stamm)}

\par\quad このケーキ、美味しそうです。\\
\par\quad Dieser Kuchen sieht lecker aus.\\

Hier hängt そう direkt an den Wortstamm (美味し-, ohne das -い) — es beschreibt einen unmittelbaren, visuellen oder sinnlichen Eindruck. Der Sprecher hat den Kuchen noch nicht probiert, aber er sieht danach aus.

\section*{Hörensagen (an die vollständige Grundform angehängt)}

\par\quad 天気予報によると、明日は雨だそうです。\\
\par\quad Laut Wettervorhersage soll es morgen regnen.\\

Hier steht そう nach der vollständigen Form (雨だ, nicht nur 雨) — es gibt wieder, was der Sprecher von jemand anderem gehört hat, nicht was er selbst wahrnimmt.

\section*{Der entscheidende Unterschied}

Beide Verwendungen sehen ähnlich aus, aber die Anschlussstelle verrät die Bedeutung:

\par\quad 面白そうです。 → Eindruck: sieht interessant aus (Stamm + そう)\\
\par\quad 面白いそうです。 → Hörensagen: soll interessant sein (Grundform + そう)\\

Ein einziges い entscheidet zwischen "ich habe den Eindruck" und "ich habe gehört, dass". Wer diese Anschlussregel kennt, kann beide Verwendungen sicher auseinanderhalten.


\chapter{たForm — Die Vergangenheitsform}

\section*{Das Bild}

た markiert, dass etwas bereits geschehen ist — die Vergangenheit. Dahinter steckt immer derselbe Gedanke: ein Vorgang liegt, aus Sicht des Sprechers im Moment des Sprechens, bereits zurück.

\section*{Nur zwei Zeitstufen}

Das Deutsche unterscheidet Vergangenheit, Gegenwart und Zukunft als drei getrennte Zeiten. Japanisch kennt grammatisch nur eine einzige Grenze: \textbf{Vergangenheit} gegenüber \textbf{Nicht-Vergangenheit}. Die Grundform eines Verbs (行く, "gehen") deckt sowohl Gegenwart als auch Zukunft ab:

\par\quad 明日、東京に行く。\\
\par\quad Ich fahre morgen nach Tokyo. (Zukunft, Grundform)\\

\par\quad いつも歩いて行く。\\
\par\quad Ich gehe immer zu Fuß. (Gegenwart/Gewohnheit, Grundform)\\

た dagegen markiert ausschließlich, dass etwas bereits geschehen ist:

\par\quad 昨日、映画を見た。\\
\par\quad Gestern habe ich einen Film gesehen.\\

Die einzige grammatische Grenze verläuft also zwischen "schon geschehen" und "noch nicht geschehen" — ob etwas gerade jetzt oder erst in der Zukunft passiert, wird von た gar nicht unterschieden und muss, falls nötig, durch Zeitangaben (明日, いつも usw.) im Satz deutlich gemacht werden.

\section*{Auch ein gerade erreichter Zustand ist schon Vergangenheit}

Jeder Zustand, der gerade erreicht wurde, liegt per Definition bereits einen Moment zurück, sobald man ihn feststellt:

\par\quad あ、財布が落ちた！\\
\par\quad Ah, die Brieftasche ist heruntergefallen!\\

Auch wenn dieser Ausruf im Moment des Bemerkens fällt, liegt das Herunterfallen selbst bereits einen winzigen Moment zurück — die Brieftasche liegt schon am Boden, wenn man es feststellt. た erfasst genau diesen Punkt: die Handlung ist bereits abgeschlossen, ihr Ergebnis ist gerade sichtbar geworden.

\par\quad 結婚しています。\\
\par\quad Ich bin verheiratet.\\

Auch hier: ein Zustand, der aus einer abgeschlossenen Handlung entstanden ist und weiterhin besteht — dieselbe Vergangenheitslogik, nur auf einen andauernden Zustand angewendet.

\textbf{Ein alltägliches Beispiel:} Wenn man gerade etwas verstanden hat, sagt man 分かりました, nicht 分かります — selbst wenn das Verstehen gerade erst, vielleicht vor einer Zehntelsekunde, geschehen ist. Im Moment, in dem man diesen Satz ausspricht, liegt der eigentliche Verstehens-Vorgang bereits hinter einem — man befindet sich schon im Zustand des Verstandenhabens, und dieser Zustand ist, wie jeder erreichte Zustand, bereits Vergangenheit.

\section*{Zur Herkunft}

た geht historisch auf eine Kombination aus て und ある zurück — wörtlich "eine Handlung ist vollzogen und der daraus entstandene Zustand besteht". Diese Verbindung zu ある erklärt, warum た von Anfang an so eng mit der Vorstellung eines bestehenden, erreichten Zustands verknüpft ist.


\chapter{てForm — Die Verkettung}

\section*{Das Bild}

て verbindet zwei Handlungen oder Zustände so, dass eine zeitliche oder inhaltliche Abfolge mitschwingt: erst das eine, dann (oder deswegen) das andere. Anders als と (Kapitel と), das zwei Elemente wirklich gleichrangig und austauschbar nebeneinanderstellt, trägt て immer eine Spur von Reihenfolge oder Zusammenhang.

\section*{Die wichtigsten Fälle}

\textbf{1. Zeitliche Abfolge}
\par\quad 朝ご飯を食べて、学校に行きます。\\
\par\quad Ich esse Frühstück und gehe (dann) zur Schule.\\

Erst das eine, dann das andere — die Reihenfolge ist festgelegt, im Gegensatz zu einer echten Aufzählung.

\textbf{2. Grund}
\par\quad 熱があって、休みました。\\
\par\quad Weil ich Fieber hatte, habe ich mich ausgeruht.\\

Hier verschiebt sich die Abfolge von rein zeitlich zu ursächlich — dasselbe Verkettungsprinzip, nur mit Betonung auf der Kausalität statt der reinen Reihenfolge.

\textbf{3. Begleitende Handlung}
\par\quad 音楽を聴いて、勉強しています。\\
\par\quad Ich lerne, während ich Musik höre.\\

Zwei Dinge, die eng miteinander verbunden ablaufen — auch hier steckt eine Verkettung dahinter, nicht zwei unabhängige, austauschbare Fakten.

\textbf{4. Höfliche Bitte}
\par\quad ちょっと待ってください。\\
\par\quad Bitte warte einen Moment.\\

てください baut direkt auf demselben Verkettungsprinzip auf: die Handlung soll sich an die Aufforderung anschließen.

\section*{Der Unterschied zu と}

と (Kapitel と) setzt zwei Elemente wirklich gleichrangig — bei einer Aufzählung mit と ist die Reihenfolge frei vertauschbar, ohne dass sich die Bedeutung ändert. Bei て dagegen steckt fast immer eine Richtung mit drin: erst dies, dann das, oder wegen diesem, jenes. Deshalb passt て besser zu Sätzen, die tatsächlich eine Abfolge oder einen Zusammenhang ausdrücken sollen, während と für reine, neutrale Aufzählung geeigneter ist.


\chapter{ばForm — Die Bedingung}

\section*{Das Bild}

ば markiert eine Bedingung: "wenn X, dann Y" — eine allgemeine, oft eher regelhafte Verknüpfung zwischen zwei Umständen.

\section*{Beispiele}

\par\quad 春になれば、桜が咲きます。\\
\par\quad Wenn der Frühling kommt, blühen die Kirschblüten.\\

Eine Bedingung, die eher wie eine allgemeine Regel klingt: sobald die eine Bedingung eintritt, folgt die andere fast automatisch.

\par\quad ボタンを押せば、ドアが開きます。\\
\par\quad Wenn man den Knopf drückt, öffnet sich die Tür.\\

\par\quad 安ければ、買います。\\
\par\quad Wenn es billig ist, kaufe ich es.\\

\section*{Der Charakter von ば}

ば klingt oft etwas allgemeiner und regelhafter als andere Bedingungsformen — passend für Sprichwörter, Naturgesetze, feste Regeln oder Anleitungen, wo eine Bedingung fast mechanisch zu einer Konsequenz führt.


\chapter{なら(ば)Form — Bedingung und Thema}

\section*{Das Bild}

なら funktioniert im Kern wie ば aus dem vorigen Kapitel, nur dass die Bedingung nicht an ein Verb, sondern an eine Feststellung ("es ist so") anknüpft. Daraus ergeben sich zwei eng verwandte Verwendungen: eine Bedingung, und — etwas unerwarteter — eine Art Themamarkierung.

\section*{Bedingung}

\par\quad 雨なら、行きません。\\
\par\quad Wenn es regnet, gehe ich nicht.\\

Dieselbe Bedingungs-Logik wie bei ば, nur angewendet auf eine bereits feststehende Situation ("es ist Regen") statt auf ein Geschehen.

\section*{Thema}

\par\quad 母なら、間もなく帰ります。\\
\par\quad Was Mutter betrifft, sie kommt bald zurück.\\

Hier wirkt なら fast wie は — es greift ein Thema heraus und stellt eine Aussage dazu in Aussicht. Diese zweite Verwendung ist keine zufällige Nebenbedeutung, sondern liegt schon in der Bedingungslogik selbst begründet: "wenn wir von X sprechen..." ist letztlich auch eine Art Bedingung, nur auf ein Gesprächsthema statt auf ein Ereignis angewendet.


\chapter{るForm — Die Möglichkeitsform}

\section*{Das Bild}

Manche Verben haben eine eigene, direkte Form, um "können" auszudrücken — nicht als Zusatzkonstruktion, sondern als eingebaute Möglichkeitsform des Verbs selbst.

\section*{Beispiele}

\par\quad 漢字が読める。\\
\par\quad Ich kann Kanji lesen.\\

読める ist die eigenständige Möglichkeitsform von 読む ("lesen") — kein Passiv, keine Höflichkeitsform, sondern ein eigenes Verb, das direkt "kann lesen" bedeutet.

\par\quad 泳げる。\\
\par\quad Ich kann schwimmen.\\

Dasselbe Muster: 泳ぐ ("schwimmen") wird zur eigenen Möglichkeitsform 泳げる.

\section*{Eine wichtige Abgrenzung}

Diese Möglichkeitsform ist nicht dasselbe wie れる/られる (das Kapitel dazu folgt als Nächstes) — beide führen letztlich zur Bedeutung "können", aber sie sind zwei unterschiedliche, unabhängig entstandene Formen. るForm entsteht direkt aus dem Verbstamm selbst und beschreibt eine Eigenschaft, die im Objekt oder in der Handlung selbst liegt — "das ist lesbar", nicht "mir wird das Lesen ermöglicht". Wer beide Formen für dasselbe hält, verwechselt zwei sprachlich getrennte Wege zum selben Bedeutungsfeld.


\chapter{れる/られる — Vier Bedeutungen, ein Ursprung}

\section*{Das Bild}

れる (bzw. られる, je nach Verbgruppe) geht auf den Gedanken zurück, dass etwas "von selbst entsteht" oder "von selbst geschieht" — ohne dass jemand es aktiv herbeiführt. Aus diesem einen Kerngedanken entwickeln sich vier auf den ersten Blick sehr unterschiedliche Verwendungen.

\section*{Die vier Verwendungen}

\textbf{1. Passiv}
\par\quad 先生に褒められました。\\
\par\quad Ich wurde vom Lehrer gelobt.\\

Die Handlung "entsteht" am Subjekt, ohne dass dieses selbst aktiv wird — sie geschieht an ihm.

\textbf{2. Spontane Empfindung (自発)}
\par\quad 昔のことが思い出される。\\
\par\quad Die alten Zeiten kommen mir in den Sinn (von selbst).\\

Hier zeigt sich die Grundidee am direktesten: eine Empfindung oder ein Gedanke entsteht von selbst, ohne bewusste Absicht.

\textbf{3. Möglichkeit}
\par\quad この水は飲まれる。\\
\par\quad Dieses Wasser ist trinkbar.\\

Eine Möglichkeit "entsteht" — etwas wird realisierbar, ohne dass jemand aktiv etwas dafür tut.

\textbf{4. Ehrerbietung}
\par\quad 先生が話される。\\
\par\quad Der Lehrer spricht. (höflich)\\

Bei einer respektierten Person wird die Handlung so dargestellt, als würde sie "von selbst" geschehen, statt sie der Person direkt und unmittelbar zuzuschreiben — eine indirekte, höfliche Distanzierung.

\section*{Warum alle vier zusammengehören}

Was Passiv, spontane Empfindung, Möglichkeit und Ehrerbietung verbindet, ist derselbe Grundgedanke: die Handlung wird nicht als aktiv vom Subjekt ausgehend dargestellt, sondern als etwas, das geschieht oder entsteht. Bei der Ehrerbietung wird diese Distanzierung sogar bewusst als Höflichkeitsmittel genutzt — man vermeidet es, die Handlung der geehrten Person zu direkt zuzuschreiben, indem man sie sprachlich "von selbst geschehen" lässt.


\chapter{せる/させる — Jemanden zum Handeln bringen}

\section*{Das Bild}

せる (bzw. させる) macht aus einem Verb eine Handlung, bei der jemand eine andere Person zu etwas veranlasst — vom starken Zwang bis zur bloßen Erlaubnis, je nachdem, wie die Konstruktion verwendet wird.

\section*{Beispiel}

\par\quad 子供に野菜を食べさせます。\\
\par\quad Ich lasse das Kind Gemüse essen. / Ich bringe das Kind dazu, Gemüse zu essen.\\

Zwischen Zwang und Erlaubnis liegt hier eine ganze Bandbreite — der Satz allein lässt beides zu, Ton und Kontext entscheiden.

\section*{Eine Skala statt eines festen Punkts}

Im Kapitel を haben wir gesehen, dass die Partikelwahl beim Kausativ diese Bandbreite sichtbar macht:

\par\quad 子供を勉強させます。 → mit を: eher Zwang, das Kind wird zum Lernen "hindurchgezwungen".\\
\par\quad 子供に勉強させます。 → mit に: eher Erlaubnis, das Kind darf lernen.\\

Kausativ und Erlaubnis sind also keine zwei getrennten Bedeutungen, sondern zwei Enden derselben Skala — starker Zwang auf der einen, lockere Duldung auf der anderen Seite. Welches Ende gemeint ist, zeigt oft gerade die Partikelwahl.

\section*{Eine höfliche Erweiterung}

\par\quad 帰らせていただきます。\\
\par\quad Ich möchte mich (mit Ihrer Erlaubnis) zurückziehen.\\

Hier wird die Kausativform mit いただく kombiniert — wörtlich "ich empfange die Erlaubnis, [zu gehen]". Das drückt besondere Bescheidenheit aus: der Sprecher bittet demütig um Erlaubnis für die eigene Handlung, statt sie einfach zu vollziehen.


\chapter{欲しい / 欲しがる — Der Wunsch nach einem Ding}

\section*{Das Bild}

欲しい drückt aus, dass man sich ein Ding wünscht — das Gegenstück zu たい (Kapitel たいForm), das sich auf eine Handlung bezieht, während 欲しい sich auf ein Ding bezieht.

\section*{Beispiel}

\par\quad 新しい靴が欲しいです。\\
\par\quad Ich möchte neue Schuhe (haben).\\

\section*{Dieselbe Grenze wie bei たい}

Genau wie たい wird 欲しい in der Regel nur für die eigene Person verwendet, oder im Gespräch für den Zuhörer:

\par\quad 何が欲しいですか。\\
\par\quad Was möchtest du (haben)?\\

Für eine dritte Person braucht es wieder die がる-Form, die wir schon von たがる kennen:

\par\quad 彼女は新しい靴が欲しがっています。\\
\par\quad Sie scheint sich neue Schuhe zu wünschen. (von außen beobachtet)\\

欲しがる beschreibt den von außen sichtbaren Wunsch einer anderen Person, nicht den direkt gewussten eigenen Wunsch — dasselbe Prinzip wie bei たい/たがる, hier auf Dinge statt auf Handlungen angewendet.


\chapter{みたいだ — Die umgangssprachliche Variante von ようだ}

\section*{Das Bild}

みたいだ funktioniert wie ようだ aus dem Kapitel よう — Vergleich, Wahrscheinlichkeit, Anschein — nur lockerer und alltagssprachlicher im Ton.

\section*{Beispiele}

\par\quad 雨が降るみたいだ。\\
\par\quad Sieht so aus, als würde es regnen.\\

\par\quad 猫みたいに寝ている。\\
\par\quad Er schläft wie eine Katze.\\

Dieselben Bedeutungen wie ようだ/ように, nur in der gesprochenen, informellen Alltagssprache üblicher.

\section*{Woher der lockere Ton kommt}

みたいだ ist tatsächlich aus 見たようだ ("es sieht aus wie gesehen") entstanden und im Lauf der Zeit lautlich zu みたいだ verschliffen worden — kein bloß ähnliches Wort, sondern eine direkte, umgangssprachliche Abwandlung von ようだ selbst. Der lockere Charakter erklärt sich also unmittelbar aus dieser Herkunft: eine im Alltag abgeschliffene, bequemere Aussprache derselben Grundidee.


\chapter{だ / です / である — Feststellung in drei Registern}

\section*{Das Bild}

だ, です und である sind im Kern dasselbe: sie stellen fest, dass etwas so ist — "X ist Y". Der Unterschied zwischen den dreien liegt nicht in der Bedeutung, sondern ausschließlich im Höflichkeitsregister.

\section*{Die drei Formen}

\par\quad これは本だ。\\
\par\quad Das ist ein Buch. (direkt, salopp)\\

\par\quad これは本です。\\
\par\quad Das ist ein Buch. (höflich, Standard)\\

\par\quad これは本である。\\
\par\quad Dies ist ein Buch. (formell, schriftsprachlich — Aufsätze, Berichte, offizielle Texte)\\

Alle drei sagen inhaltlich exakt dasselbe — nur der Ton unterscheidet sich.

\section*{Warum だ direkter wirkt, als man zunächst denkt}

だ hat den Ruf, besonders direkt oder abrupt zu klingen, fast wie ein Punkt hinter der Aussage: "Das ist so. Fertig." Ein naheliegender, aber falscher Gedanke wäre: だ klingt härter, weil man es nicht verneinen kann. Das stimmt so nicht — だ lässt sich sehr wohl verneinen (これは本ではない/じゃない), genau wie です (これは本ではありません). Beide brauchen dafür ein eigenes Ersatzwort (ではない/じゃない), keins der beiden hat eine eigene eingebaute Verneinungsform.

Der eigentliche Grund für den direkten Klang liegt woanders: だ wird in natürlicher Rede sehr häufig einfach weggelassen, um den Satz weicher klingen zu lassen (本。 statt 本だ。) — です dagegen fällt praktisch nie weg, weil es der höfliche Standardabschluss ist. Diese Auslassbarkeit, zusammen mit dem allgemein als burschikos-direkt empfundenen Ton, prägt den Eindruck, dass だ "härter" ist — nicht die Verneinung.


\chapter{Höflichkeitsformen — Wer wird gehoben, wer senkt sich}

\section*{Das Bild}

Japanische Höflichkeit funktioniert nicht über einen einzigen Höflichkeitsgrad, sondern über die Frage: wen hebt man an, und wer senkt sich ab? Es gibt drei grundlegend verschiedene Bewegungsrichtungen, aus denen sich die Höflichkeitsformen zusammensetzen.

\section*{Die drei Bewegungsrichtungen}

\textbf{1. Den anderen erhöhen (尊敬語)}
\par\quad 先生がいらっしゃいます。\\
\par\quad Der Lehrer kommt. (respektvoll, weil es sich um eine höhergestellte Person handelt)\\

Die Handlung des Gegenübers oder einer Respektsperson wird sprachlich "angehoben" — eigene, besondere Verbformen ersetzen die neutralen (いらっしゃる statt 来る/行く/いる).

\textbf{2. Sich selbst absenken (謙譲語)}
\par\quad 先生のところに伺います。\\
\par\quad Ich werde zum Lehrer gehen. (bescheiden, weil man die eigene Handlung gegenüber der Respektsperson zurücknimmt)\\

Hier senkt sich der Sprecher sprachlich ab, um dadurch indirekt die andere Person zu erhöhen — dieselbe Grundidee wie bei 尊敬語, nur aus der umgekehrten Richtung angegangen.

\textbf{3. Allgemeine Höflichkeit (丁寧語)}
\par\quad 明日、行きます。\\
\par\quad Ich gehe morgen. (einfach höflich, ohne jemanden besonders zu erhöhen oder sich selbst zu senken)\\

Das kennen wir bereits aus dem Kapitel zur ますForm — höflicher Umgangston gegenüber dem Zuhörer, unabhängig davon, wer in der Handlung selbst vorkommt.

\section*{Eine Feinheit bei der Selbstabsenkung}

Bescheidene Sprache (謙譲語) gibt es in zwei Varianten, die sich in einem wichtigen Punkt unterscheiden:

\textbf{Gerichtete Bescheidenheit} — die Handlung bezieht sich konkret auf eine bestimmte Person, die dadurch erhöht wird:
\par\quad 先生に伺います。\\
\par\quad Ich werde den Lehrer fragen. (senkt sich explizit vor dem Lehrer)\\

\textbf{Ungerichtete Bescheidenheit} — ein allgemein formeller, bescheidener Sprachstil, ohne dass eine konkrete Person angesprochen oder erhöht wird:
\par\quad 明日、参ります。\\
\par\quad Ich werde morgen kommen. (formeller Sprechstil, kein bestimmtes Gegenüber wird erhöht)\\

Der Unterschied lässt sich als Bild festhalten: die erste Variante ist wie eine Verbeugung vor einer bestimmten Person, die zweite eher wie ein Anzug, den man bei jedem formellen Anlass trägt — unabhängig davon, wer konkret anwesend ist.


\chapter{Adjektive — Zwei ganz unterschiedliche Familien}

\section*{Das Bild}

Japanisch hat zwei völlig verschiedene Arten von "Adjektiven", die sich äußerlich ähneln, aber grammatisch anders funktionieren: い-Adjektive, die sich selbst wie Verben konjugieren, und な-Adjektive, die sich eher wie Nomen verhalten.

\section*{い-Adjektive}

\par\quad 高い山\\
\par\quad ein hoher Berg\\

\par\quad この山は高い。\\
\par\quad Dieser Berg ist hoch.\\

い-Adjektive können direkt vor einem Nomen stehen (attributiv) oder direkt einen Satz beenden (prädikativ), ohne dass es dazwischen eine Kopula (だ/です) braucht. Sie tragen ihre eigene Konjugation in sich (高かった, "war hoch"; 高くない, "ist nicht hoch").

\section*{な-Adjektive}

\par\quad 静かな部屋\\
\par\quad ein ruhiges Zimmer\\

\par\quad この部屋は静かです。\\
\par\quad Dieses Zimmer ist ruhig.\\

な-Adjektive brauchen vor einem Nomen ein angehängtes な, und im Satzabschluss die Kopula だ/です — genau wie ein gewöhnliches Nomen (本の部屋 hätte dieselbe Struktur wie 静かな部屋, nur mit の statt な).

\section*{Warum な-Adjektive eigentlich Nomen sind}

Diese Ähnlichkeit zu gewöhnlichen Nomen ist kein Zufall. In der japanischen Grammatikforschung gibt es die anerkannte Position, dass な-Adjektive gar keine eigene Wortart bilden, sondern schlicht Nomen sind, die mit だ verbunden werden — 静かな ist demnach eigentlich das Nomen 静か plus eine Form von だ, genau wie ein normales Nomen. Diese Sichtweise wird durch die Konjugation selbst gestützt: die Flexionsformen eines な-Adjektivs (だろ, だっ/で/に, だ, な, なら) sind identisch mit der Flexion der Kopula だ selbst, die wir bereits aus dem Kapitel だ/です/である kennen.

Für die Praxis heißt das: な-Adjektive lassen sich am einfachsten behandeln, indem man sie sich als Nomen mit besonders enger Beziehung zu だ vorstellt — nicht als eigenständige, verbähnliche Wortart wie die い-Adjektive.


\chapter{Nomen und vs-Nomen — Wenn ein Nomen zum Verb wird}

\section*{Das Bild}

Viele japanische Nomen — meist chinesischen Ursprungs — lassen sich direkt mit する ("tun") zu einem Verb kombinieren. Statt für jede Handlung ein eigenes Verb zu lernen, reicht oft: Nomen + する.

\section*{Der wichtigste Fall}

\par\quad 運動をします。\\
\par\quad Ich mache Sport.\\

\par\quad 運動します。\\
\par\quad Ich mache Sport. (gleichbedeutend)\\

Das を ist hier optional — solange kein anderes Objekt im Satz auftaucht, kann man es weglassen, ohne dass sich die Bedeutung ändert.

Diese Optionalität gilt allerdings nur für das nackte Nomen. Sobald ein Adjektiv das Nomen näher beschreibt, wird を verpflichtend:

\par\quad 難しい運動をします。\\
\par\quad Ich mache anstrengenden Sport.\\

\par\quad ×難しい運動します。\\
\par\quad (nicht korrekt)\\

Sobald also ein Adjektiv davorsteht, muss を stehen bleiben — nur beim unveränderten, alleinstehenden Nomen kann es wegfallen.

\section*{Eine Erweiterung: Wahrnehmung und Eindruck}

\par\quad いいにおいがする。\\
\par\quad Es riecht gut.\\

\par\quad 変な音がする。\\
\par\quad Es klingt seltsam.\\

Hier funktioniert する anders — nicht als "eine Handlung ausführen", sondern als "eine Wahrnehmung entsteht". Trotzdem bleibt das Grundprinzip erhalten: ein Nomen (におい, "Geruch") wird über する zu einem vollständigen Ausdruck.

\section*{Eine produktive, lebendige Konstruktion}

Diese Nomen+する-Bildung ist im modernen Japanisch weiterhin sehr produktiv — sie wird auch auf Nomen angewendet, die traditionell nicht dafür gedacht waren:

\par\quad 青春する。\\
\par\quad Die Jugend in vollen Zügen erleben. (wörtlich: "Jugend tun")\\

Solche kreativen Bildungen zeigen, wie flexibel dieses einfache Baumuster ist: fast jedes passende Nomen lässt sich mit する zu einer eigenen Handlung machen.


\chapter{Adverbien — Direkt, abgeleitet, oder als Nomen verkleidet}

\section*{Das Bild}

Japanische Adverbien kommen in drei Formen: eigenständige Adverbien, aus Adjektiven abgeleitete Adverbien, und eine dritte, weniger offensichtliche Gruppe — Nomen, die adverbiale Funktionen übernehmen.

\section*{Eigenständige Adverbien}

\par\quad とても嬉しいです。\\
\par\quad Ich bin sehr glücklich.\\

Wörter wie とても ("sehr") funktionieren als eigenständige, unveränderliche Adverbien.

\section*{Von Adjektiven abgeleitete Adverbien}

\par\quad 静かに話します。\\
\par\quad Ich spreche leise.\\

Aus dem な-Adjektiv 静か wird über に ein Adverb — dieselbe Logik, die wir schon bei に als Bezugspunkt-Partikel gesehen haben, angewendet auf den Übergang von einer Eigenschaft in eine Art und Weise.

\section*{Nomen mit adverbialer Funktion}

\par\quad 食べる時、手を洗います。\\
\par\quad Wenn ich esse, wasche ich mir die Hände. (wörtlich: "zur Zeit des Essens")\\

時 ("Zeit") wirkt hier wie ein Adverb, ist aber eigentlich ein Nomen mit stark abgeschwächter Eigenbedeutung — es dient hauptsächlich dazu, den vorangehenden Satz grammatisch anschlussfähig zu machen, ähnlich wie es の im Kapitel の tut. Wörter wie 後 ("danach"), 間 ("während"), 場合 ("Fall, Situation") funktionieren nach demselben Prinzip: sie sind reichere Varianten desselben Grundmechanismus, mit einer eigenen zeitlichen oder logischen Bedeutung, aber strukturell demselben Trick wie の — ein Nomen erlaubt es, einen Satz an eine Partikel anzuschließen.


\chapter{あげる / くれる / もらう — Wessen Augen sehen die Handlung}

\section*{Das Bild}

Bei Geben und Bekommen wählt man im Japanischen das Verb nicht nach Rang oder Höflichkeit allein, sondern zuerst danach, aus wessen Perspektive die Handlung erzählt wird — und ob der Sprecher sich überhaupt in diese Perspektive hineinversetzen kann.

\section*{Die drei Grundverben}

\textbf{あげる — aus Geber- oder neutraler Sicht}
\par\quad 私は友達に本をあげました。\\
\par\quad Ich habe meinem Freund ein Buch gegeben.\\

あげる erzählt aus der Perspektive dessen, der gibt — oder bleibt neutral, ohne sich stark in eine der beiden Seiten hineinzuversetzen.

\textbf{くれる — nur aus Empfänger-Sicht}
\par\quad 友達が私に本をくれました。\\
\par\quad Mein Freund hat mir ein Buch gegeben.\\

くれる verlangt zwingend, dass der Sprecher die Perspektive des Empfängers einnimmt — er muss sich in dessen Position hineinversetzen können. Deshalb wird くれる praktisch immer verwendet, wenn der Sprecher selbst (oder jemand ihm sehr Nahestehendes) der Empfänger ist: man kann sich in die eigene Position natürlich am leichtesten hineinversetzen.

\textbf{もらう — Empfänger-Perspektive, unabhängig von der Richtung}
\par\quad 私は友達に本をもらいました。\\
\par\quad Ich habe von meinem Freund ein Buch bekommen.\\

もらう blickt immer vom Empfänger aus auf das Geschehen, unabhängig davon, wer gibt.

\section*{Warum くれる so eng an "die eigene Gruppe" gebunden wirkt}

Weil くれる zwingend die Fähigkeit verlangt, sich in den Empfänger hineinzuversetzen, wird es in der Praxis fast nur verwendet, wenn der Empfänger der Sprecher selbst oder jemand aus seiner unmittelbaren Nähe ist — das ist keine feste Regel über Gruppenzugehörigkeit, sondern eine natürliche Folge der Blickwinkel-Anforderung: Menschen versetzen sich am leichtesten in die eigene Position oder die von Menschen, die ihnen nahestehen.

あげる dagegen ist flexibler und kann auch verwendet werden, wenn weder Geber noch Empfänger der Sprecher selbst sind — solange die Erzählung neutral bleibt.

\section*{Die zweite Achse: Rang}

Innerhalb dieser durch die Perspektive festgelegten Grundrichtung gibt es noch eine zweite Ebene: das konkrete Verb ändert sich je nach Rangverhältnis zwischen Geber und Empfänger.

\par\quad 先生に本を差し上げました。\\
\par\quad Ich habe dem Lehrer (ehrerbietig) ein Buch gegeben.\\

\par\quad 猫にえさをやりました。\\
\par\quad Ich habe der Katze Futter gegeben. (neutral bis leicht herablassend, üblich bei Tieren, Pflanzen)\\

差し上げる ist die respektvolle Variante von あげる bei höhergestellten Empfängern, やる eine neutrale bis leicht herablassende Variante bei Tieren oder Pflanzen. Dieselbe Logik gilt spiegelbildlich bei くれる (下さる als respektvolle Variante) und もらう (いただく als bescheidene Variante). Der Rang bestimmt also nicht die Grundrichtung selbst, sondern nur, welches konkrete Wort innerhalb dieser Richtung passend ist.

\end{document}
